\section{Conclusion}
\label{sec:conclusion}

We set out to build a thermal wildlife counting system and found that the accuracy a
practitioner would try to buy with a better detector is being lost somewhere else. By
decomposing counting error into detection, association and confirmation, we showed that on
held-out video the pipeline finds 98.8\% of animals, gives 88.0\% a distinct track, and
counts 62.7\% --- and that four independent improvements to the first two stages each made
the third worse.

The evidence points to supervision scarcity rather than model class as the reason
confirmation resists learning. Pool C offers 207 positive tracks against 6,801 negatives and
pool G 219 against 27,460, and the positives are not cleanly separable: duplicate tracks,
fragments of an animal some other track already covers, sit \emph{between} the positives and
the false candidates in every feature we measured, at median confidence 0.42 against 0.72
for primaries and 0.20 for false candidates. A confirmer therefore receives contradictory
supervision on exactly the examples that distinguish counting from detection, which is
consistent with error growing monotonically with capacity and with the fact that enlarging
the pool --- multiplying negatives without multiplying positives --- makes every method worse.
The capacity trend extrapolates to a crossover in the low thousands of positive tracks, an
order of magnitude beyond this corpus and a concrete target for the next dataset.

The practical guidance is inverted from the field's default. For a counting application at
this data scale: do not select a detector by its detection metric, because the counting
criterion re-ranks the models; do not evaluate on video the detector trained on, which
inflates coverage by 23 points; do not assume a larger candidate pool is a better one; and
do not expect a learned confirmation stage to beat a hand-tuned rule until the corpus is
substantially larger than 236 individuals.

Several limits bound these claims. The corpus is 236 animals from four sites in one region
and one winter, one site contributes 56\% of them, and per-site figures for the smallest
site (15 animals) are correspondingly noisy; all results are single-species. Our headline
counting metric is a per-video aggregate, and the stricter identity-matched figure is 10
points lower. No individual animal appears in more than one video, so the re-identification
limit of \cref{sec:reid} is measured within video and cross-session identification could not
be evaluated at all.

We release the dataset, per-individual annotations, evaluation harness and decomposition
tooling. The gap between 88\% reachable and 63\% counted is the most useful thing we can
hand to the next group working on this problem: it is large, precisely localized, and
entirely open.
